\documentclass{homework}
% \usepackage{graphicx} % Required for inserting images
\usepackage{amsmath}
\usepackage{gensymb}
\usepackage{subfig}
\usepackage{float}
\usepackage{amsthm}

\class{ECEN 5478}
\title{Homework 3}
\author{Sean Jennings}
\date{October 3, 2023}

\begin{document}
\maketitle

\section{Problem 1}

\newcommand{\convexarg}{\lambda x + (1 - \lambda y)}
\newcommand{\convexineq}[2]{
    #1(\lambda x + (1 - \lambda y)) #2 \leq \lambda #1 (x)
        + (1 - \lambda) #1 (y)
}
\newcommand{\strongconvexterm}[1]{
    \frac{#1}{2} \lambda (1 - \lambda) \norm{y - x}^2
}

\begin{letterenum}
    \item By the definition of convexity and $\mu$-strong convexity,
        we have:
        \begin{align}
            \convexineq{f}{&} - \strongconvexterm{\mu}
                \label{eq:f-strong-convex}\\
            \convexineq{g}{&} \label{eq:g-convex}
        \end{align}
        Adding eq (\ref{eq:f-strong-convex}) and (\ref{eq:g-convex})
        yields
        \begin{align*}
            (f+g) (\convexarg) &= f(\convexarg) + g(\convexarg) \\
                &\leq \lambda f(x) + (1 - \lambda)f(y) - \strongconvexterm{\mu} \\
                &\quad + \lambda g(x) + (1 - \lambda)g(y) \\
                &= \lambda (f+g)(x) + (1 - \lambda)(f+g)(y) - \strongconvexterm{\mu}
        \end{align*}
        so that $f + g$ is $\mu$-strongly convex, with strong convexity
        coefficient equivalent to that of $f$.

    \item Again by definition of strong convexity, we have:
    \begin{equation}
        \convexineq{g} - \strongconvexterm{\nu}.
        \label{eq:g-strongly-convex}
    \end{equation}
    Following a similar argument to part a), instead adding equations
    (\ref{eq:f-strong-convex}) and (\ref{eq:g-strongly-convex}), we have
    \begin{align*}
        (f + g) (\convexarg)
            &\leq \lambda f(x) + (1 - \lambda)f(y) - \strongconvexterm{\mu} \\
            &\quad + \lambda g(x) + (1 - \lambda)g(y) - \strongconvexterm{\nu} \\
            &= \lambda (f + g)(x) + (1 - \lambda) (f + g)(y) - \strongconvexterm{\mu + \nu}
    \end{align*}
    thus proofing $f + g$ is $(\mu + \nu)$-strongly convex.
\end{letterenum}

\section{Problem 2}
For $f$ defined
\begin{equation}
    f(x) = \frac{1}{2} x^TQx - a \norm{x}^2,
\end{equation}
$f$ is twice continuously differentiable with gradient
\begin{equation}
    \nabla f(x) = Qx - 2ax
\end{equation}
and hessian
\begin{equation}
    \nabla^2 f(x) = Q - 2aI.
\end{equation}
where $I\in\R^{n\times n}$ is the identity matrix.

We note a preliminary fact: if $A\in\R^{n\times n}$ has eigenvectors
$v_i$ for $i=1,\dots,n$ and corresponding eigenvalues $\lambda_i$,
then for any $c \in \R$, each eigenvector $v_i$ of $A$ is also
an eigenvector of $A - cI$ with eigenvalue $\lambda_i - c$.

\begin{proof}
    \begin{equation}
        (A - cI) v_i = Av_i - cv_i = \lambda v_i - cv_i = (\lambda - c) v_i
    \end{equation}
\end{proof}

\begin{letterenum}
    \item Since $f$ is twice continuously differentiable, and the
        domain in question is $\R^n$, we can check second order
        convexity conditions. The hessian must be positive semi-definite:
        \begin{equation}
            \nabla^2 f(x) = Q - 2aI \succeq 0
        \end{equation}
        By the eigenvector characterization theorem, we have that the minimum
        eigenvalue of $Q - 2aI$ must be non-negative. Since the eigenvalues
        of $Q-2aI$ are related to those of $Q$ by the fact proved above, we have
        \begin{equation}
            \lambda_{min} \geq 2a
        \end{equation}
        where $\lambda_{min}$ is the minimum eigenvalue of $Q$.

    \item Consider the function
    \begin{equation}
        g(x) := f(x) - \frac{\mu}{2}\norm{x}^2
    \end{equation}
    for some $\mu\in\R$ with $\mu > 0$.
    
    By the theorem relating strong convexity to the definition of
    $g$, we have $f$ is $\mu$-strongly convex if and only if $g$
    is convex. It is clear that $g$ is also twice continuously
    differentiable with hessian
    \begin{equation}
        \nabla^2 g(x) = Q - (2a + \mu) I.
    \end{equation}
    and thus $g$ is convex iff $Q - (2a + \mu)I \succeq 0$. Again by eigenvalue
    characterization, $g$ is convex iff $\lambda_{min} - (2a + \mu) \geq 0$,
    where $\lambda_{min}$ is defined as above, which implies
    $\lambda_{min} - 2a \geq \mu > 0$.

    If $\lambda_{min} > 2a$, we can select
    $0 < \mu \leq \lambda_{min} - 2a$ such that $g$ is convex.
    If $\lambda_{min} \leq 2a$, no such $\mu$ exists.
    Conversely, if $g$ is convex, then $\lambda - 2a > 0$.

    Therefore, $Q - 2aI \succ 0$. 
    is a sufficient and necessary condition for $f$ being
    $\mu$-strongly convex. In terms of eigenvalues, this condition
    if equivalent to $\lambda_{min} > 2a$.

    We note that for $f$ to be strongly convex, the preceding
    inequality on the minimum eigenvalue of $Q$ from part a) becomes strict. Equivalently,
    $Q$ must be positive definite, rather than positive semi-definite.
\end{letterenum}

\section{Problem 3}
\begin{letterenum}
    \item Since $f$ is convex, by the gradient inequality, we have
    \begin{equation}
        f(x) + \nabla(x)^T(y-x) \leq f(y)
    \end{equation}
    for all $x,y\in\X$. Put $y=x^*$.
    \begin{align*}
        f(x) - f(x^*) \leq \nabla f(x)^T (x - x^*) \\
            &\leq \norm{\nabla f(x)} \norm{x - x^*} \\
            &\leq G(x - x^*)
    \end{align*}

    \item Define the operator $F: \R^n \to \R^n$,
    $F := \mathcal{I} - \alpha \nabla f$ so that
    $x_{k+1} = (\text{proj}_{\X} \circ F) (x_k)$. We assume for simplicity
    that $f$ is twice-continuously differentiable. We then see that the
    gradient of $F$ is:
    \begin{equation}
        \nabla F = I - \alpha \nabla^2 f
    \end{equation}
    Since $f$ is twice-continouusly differentiable, $L$-smooth and $\mu$-strongly
    convex, its minimum and maximum eigenvalues are given by
    \begin{equation}
        \lambda_{max}(\nabla^2 f) = L \qquad \lambda_{min}(\nabla^2 f) = \mu
    \end{equation}
    Then for all $x \in \R^n$ we have the following ordering
    \begin{equation}
        (1 - \alpha L) I \preceq \nabla F(x) \preceq (1 - \alpha \mu) I
    \end{equation}
    Thus, the norm of $\nabla F$ has the following bound:
    \begin{equation}
        \norm{\nabla F} \leq \rho
    \end{equation}
    where $\rho := \text{max}\{ |1 - \alpha L|, | 1 - \alpha \mu | \}$.
    
    Since its gradient is bounded by $\rho$, $F$ is Lipschitz with constant $\rho$.

    Since the projection operator is nonexpansive, the composition
    $\text{proj}_{\X} \circ F$ is also Lipschitz with constant $\rho$. Thus,
    putting $x = x_k$ into the result from the previous section, and 
    iteratively appling the Lipschitz property, we find the desired bound
    \begin{align*}
        f(x_k) - f(x^*) &\leq G \norm{x_k - x^*} \\
            &= G \norm{F(x_{k-1}) - F(x^*)} \\
            &\leq G \rho \norm{x_{k-1} - x^*} \\
            &\leq G \rho^k \norm{x_{k-1} - x^*}
    \end{align*}
\end{letterenum}

\section{Problem 4}
We rewrite the regularized least square's cost function in matrix form with
$a_i \in R^n$ as the $i$th row of the matrix $A \in R^{N \times n}$ and 
$y_i \in \R$ as the $i$th entry of $y \in \R^n$.

\begin{equation}
    f(x) = \norm{Ax - y}^2 + \frac{\lambda}{2} \norm{x - \bar{x}}^2
\end{equation}
The gradient and hessian are given as follows
\begin{align}
    \nabla f(x) &= 2A^T(Ax - y) + \lambda(x - \bar{x}) \label{eq:gradient} \\
    \nabla^2 f(x) &= 2A^TA + \lambda I \label{eq:hessian}
\end{align}

\begin{letterenum}
    \item For any $x,z \in \R^n$, we have
    \begin{align*}
        \norm{\nabla f(x) - \nabla f(z)} &=
            \norm{2A^T(Ax - y) + \lambda (x - \bar{x}) \\
        &\qquad - 2A^T(Az - y) - \lambda (z - \bar{x})} \\
        &= \norm{2A^TA (x - z) + \lambda (x - z)} \\
        &\leq \norm{2A^TA + \lambda I} \norm{x - z}.
    \end{align*}
    So it is $L$-smooth with minimum smoothness constant
    $L = 2\lambda_{max}(A^TA) + \lambda$. Since $A^TA \succeq 0$ and $\lambda > 0$,
    we have $L > 0$.

    \item To check strong convexity, we look at the minimum eigenvalue of the hessian
    from equation (\ref{eq:hessian}).
    \begin{align}
        \mu = \lambda_{min}(\nabla^2 f(x)) &= \lambda_{min}(A^TA + \lambda I) = \lambda_{min}(A^TA) + \lambda
            \label{eq:min-eigenvalue} \\
            &=  \lambda
    \end{align}
    where we used the fact that since $N < n$, $A^TA$ is not full rank and thus
    has a zero eigenvalue. Thus, it is $\mu$-strongly convex with $\mu=\lambda$.
    The regularization term therefore ensures strong convexity even when
    there are insufficient data points.

    \item Assuming that $A$ is full column rank, $N > n$ implies that $A^TA$
    is invertible and is thus positive definite. It is still $\mu$-strongly
    convex, with $\mu$ given by equation (\ref{eq:min-eigenvalue}), but now
     $\mu > \lambda$ since $\lambda_{min}(A^TA) > 0$.
\end{letterenum}



% \section{Appendix}

% Suppose $f: \R^n \to \R^n$ is Lipschitz with constant $L > 0$ and is continuously
% differentiable. Then,
% \begin{equation}
%     \norm{\nabla f(x)} \leq L
% \end{equation}
% for all $x \in R^n$.

% \begin{proof}
%     Take $x,h \in R^n$, with $h$ such that
%     $\nabla f(x)^T h = \norm{\nabla f(x)} \norm{h}$.  By the Lipschitz
%     property of $f$, we have
%     \begin{equation}
%         \norm{f(x + h) - f(x)} \leq L \norm{h}
%     \end{equation}

    
% \end{proof}




\end{document}
